\section*{Part II: Price Mathematics Without False Precision}

\subsection*{8. Price, Quantity, Marginal Cost, and Contribution}

In a simplified 1-product model with unit price $P$, unit incremental cost $c$, quantity demanded $Q(P)$, and fixed period cost $F$, period profit is
\begin{equation}
\Pi(P) = (P - c)Q(P) - F.
\end{equation}
This clarifies why revenue is not profit. A lower price may increase volume while reducing contribution. A higher price may improve per-customer economics but reduce quantity. When price changes buyer mix, unit costs and acquisition costs may also change, and the simple expression is incomplete. Operating costs sometimes jump when a 2nd reviewer, production line, or support team becomes necessary.

The useful founder habit is to test a \emph{joint} forecast: paid quantity, realized price, variable costs, acquisition cost, capacity, quality, and cash collection. Optimizing any 1 separately can make the whole worse.

\subsection*{9. The Price-Cut Break-Even Calculation}

Suppose 100 units sell at ₹1,000, with ₹400 of incremental cost each. Contribution is $100 \times ₹600 = ₹60,000$. Cut price to ₹800 and contribution per unit becomes ₹400. The firm must now sell 150 units to preserve ₹60,000 before additional costs: a \textbf{50 percent volume increase}. A price increase from ₹1,000 to ₹1,200 yields ₹800 contribution per unit; the firm could fall to 75 units and still preserve contribution.

In general, the required new volume for equal contribution, holding unit cost and fixed commitments constant, is
\begin{equation}
Q_1^* = Q_0 \left( \frac{P_0 - c}{P_1 - c} \right).
\end{equation}
That is a hurdle, not a demand forecast. If the founder claims a 20 percent discount will increase demand 70 percent, the claim needs actual evidence. If the discount also increases support load or customer acquisition cost (CAC), the hurdle rises. If it fills otherwise idle capacity with no quality deterioration, the decision may be sensible even at lower unit contribution.

\subsection*{10. Elasticity and the Lerner Relation}

Price elasticity of demand measures responsiveness: $\epsilon = (dQ/dP)(P/Q)$. In a stylized differentiable monopoly model with constant marginal cost and no other complications, profit optimization implies
\begin{equation}
\frac{P - c}{P} = -\frac{1}{\epsilon}.
\end{equation}
This is the Lerner relation. It illustrates why the optimal markup depends on how readily buyers substitute away, not just the seller's accounting cost. It is not a general-purpose startup pricing formula. Real demand may be kinked, capacity-limited, heterogeneous, and dynamic; reputation, repeat purchases, and multiproduct interactions can matter. The relevant market may be competitive rather than monopolistic, and many customer responses cannot be represented by 1 stable elasticity.

An important operational insight is that the revenue-maximizing price and contribution-maximizing price need not match. Price changes must be evaluated on margins and total economic consequences, not turnover alone.

\subsection*{11. Break-Even Enrollment Is Only the Beginning}

For a hypothetical fixed-price cohort with fixed cost $F$, price $P$, and incremental cost per learner $c$, break-even quantity is
\begin{equation}
Q_{\text{BE}} = \left\lceil \frac{F}{P - c} \right\rceil, \quad \text{provided } P > c.
\end{equation}
Suppose $F = ₹2,20,000$ and $c = ₹7,000$:
\begin{itemize}
  \item At ₹20,000: break-even is 17 learners ($\lceil ₹2,20,000 / ₹13,000 \rceil = 17$).
  \item At ₹40,000: break-even is 7 learners ($\lceil ₹2,20,000 / ₹33,000 \rceil = 7$).
  \item At ₹70,000: break-even is 4 learners ($\lceil ₹2,20,000 / ₹63,000 \rceil = 4$).
\end{itemize}
These thresholds are mathematically correct but commercially incomplete. They do not tell us whether enough qualified buyers exist, whether the institution can deliver the service, or whether higher-paying students demand more individualized support.

The more informative table considers demand at each price, cohort capacity, acquisition cost, learner success, and normalized founder labor. The apparent ``optimal'' tuition will change when those assumptions change. Founders should prefer a range of scenarios to a spreadsheet that makes invented conversion rates look authoritative.
